% !TEX program = lualatex
\documentclass[10pt,a4paper]{article}
\usepackage[T1]{fontenc}
\usepackage[french]{babel}
\usepackage[left=1cm,right=1cm,top=1.3cm,bottom=1.3cm]{geometry}
\usepackage{amsmath}
\usepackage{amssymb}
\usepackage{array}
\usepackage{graphicx}
\usepackage{xcolor}
\usepackage{titlesec}
\usepackage{enumitem}
\usepackage{multicol}
\raggedcolumns
\usepackage{pifont}
\usepackage{tikz}
\usetikzlibrary{calc,arrows.meta,shapes.misc}
\tikzset{ptcross/.style={cross out,draw,thick,minimum size=3pt,inner sep=0pt}}
% Badges « Exercice n » en cursive Playwrite FR Moderne (police du site),
% chargée à 85 % de la taille courante (XeLaTeX ou LuaLaTeX ; gras sous pdfLaTeX).
\makeatletter
\ifdefined\XeTeXrevision\def\ExoFont@ok{}\fi
\ifdefined\directlua\def\ExoFont@ok{}\fi
\ifdefined\ExoFont@ok
\newcommand{\ExoFont}{\font\ExoFont@f="[PlaywriteFRModerne.ttf]" at \dimexpr\f@size pt*85/100\relax\ExoFont@f}
\else
\newcommand{\ExoFont}{\bfseries}
\fi
\makeatother
% \ExoTexte{...} : texte du badge en cursive, avec la hauteur d'une majuscule de la
% police courante (la cursive est plus haute et ferait grandir chaque badge).
\newcommand{\ExoTexte}[1]{\raisebox{0pt}[\fontcharht\font`E][0pt]{\ExoFont #1}}
\usepackage{tcolorbox}
\tcbuselibrary{skins}
\usepackage[hidelinks]{hyperref}
\definecolor{calcblue}{HTML}{1C7ED6}
\definecolor{skyblue}{HTML}{339AF0}
\definecolor{objorange}{HTML}{F76707}
\definecolor{repred}{HTML}{E03131}
\newcommand{\Vj}[1]{\overrightarrow{\text{#1}}}
\newcommand{\Pt}[1]{\textsf{\footnotesize #1}}
\renewcommand{\familydefault}{\sfdefault}
\DeclareSymbolFont{sfnumbers}{OT1}{cmss}{m}{n}
\SetSymbolFont{sfnumbers}{bold}{OT1}{cmss}{bx}{n}
\DeclareMathSymbol{0}{\mathalpha}{sfnumbers}{`0}
\DeclareMathSymbol{1}{\mathalpha}{sfnumbers}{`1}
\DeclareMathSymbol{2}{\mathalpha}{sfnumbers}{`2}
\DeclareMathSymbol{3}{\mathalpha}{sfnumbers}{`3}
\DeclareMathSymbol{4}{\mathalpha}{sfnumbers}{`4}
\DeclareMathSymbol{5}{\mathalpha}{sfnumbers}{`5}
\DeclareMathSymbol{6}{\mathalpha}{sfnumbers}{`6}
\DeclareMathSymbol{7}{\mathalpha}{sfnumbers}{`7}
\DeclareMathSymbol{8}{\mathalpha}{sfnumbers}{`8}
\DeclareMathSymbol{9}{\mathalpha}{sfnumbers}{`9}
\setlength{\parindent}{0pt}
\setlength{\columnsep}{18pt}
\graphicspath{{./}}
\newcommand*{\textcircledd}[1]{\begin{tikzpicture}[baseline=(letter.base)]\node[draw,circle,inner sep=1pt] (letter) {#1};\end{tikzpicture}}
\newcommand{\numeroCerne}[1]{\protect\textcircledd{\bfseries\scriptsize #1}}
\setlist[enumerate,1]{label=\numeroCerne{\arabic*}, leftmargin=*, labelsep=5pt, itemsep=2pt, topsep=1pt, partopsep=0pt, parsep=0pt, beginpenalty=10000}
\setlist[itemize,1]{beginpenalty=10000}
\setlist[enumerate,2]{label=(\alph*), leftmargin=*, labelsep=4pt, itemsep=1pt, topsep=1pt, partopsep=0pt, parsep=0pt}
\setlength{\abovedisplayskip}{3pt}\setlength{\belowdisplayskip}{3pt}\setlength{\abovedisplayshortskip}{3pt}\setlength{\belowdisplayshortskip}{3pt}

% Ligne de corrig\'e : num\'ero d'exercice en pastille bleu ciel (m\^eme
% couleur que les cartes d'exercice de la feuille), suivi de la r\'eponse.
\newcommand{\Corr}[2]{%
\par\vspace{6pt}
\noindent\tcbox[on line, colback=skyblue, colframe=skyblue, boxsep=1pt, left=4pt, right=4pt, top=1pt, bottom=1pt, arc=2pt]{\textcolor{white}{\ExoTexte{Exercice~#1}}}\par\nobreak\vspace{3pt}\nobreak
\noindent\ignorespaces #2
\par
}

% Bandeau d'objectif, comme sur la feuille d'exercices.
\newcommand{\Objectif}[2]{%
\par\vspace{8pt}
\noindent\begin{minipage}{\linewidth}
{\large\color{objorange}\bfseries Objectif #1~:}\\{\large\color{objorange}#2}
\end{minipage}\par\nobreak
\vspace{2pt}\nobreak
}

% \'El\'ements construits (points plac\'es, vecteurs trac\'es) : en rouge.
\newcommand{\Rep}[1]{{\color{repred}#1}}

\author{2nde --- Lyc\'ee Lebrun Coutances}
\date{}
\title{\textbf{S\'equence 2 --- Feuille d'exercices n\textdegree{}3\,\textsuperscript{3/4} : Vecteurs et d\'emonstrations}\\[4pt]\large Corrig\'e}
\begin{document}
\maketitle
\vspace{-18pt}

\begin{multicols}{2}

\Objectif{5}{Produit d'un vecteur par un nombre r\'eel et colin\'earit\'e de deux vecteurs}

\Corr{1}{
\begin{minipage}{\linewidth}
\begin{enumerate}
\item $\Vj{AC}$ correspond \`a 1 carreau \`a droite et 2 en haut, donc $\Vj{BM}=2\Vj{AC}$ \`a 2 carreaux \`a droite et 4 en haut \`a partir de $B$.\\
$\Vj{BC}$ correspond \`a 2 carreaux \`a gauche et 2 en haut, donc $\Vj{AN}=\dfrac12\Vj{BC}$ \`a 1 carreau \`a gauche et 1 en haut \`a partir de $A$.
\begin{center}
\begin{tikzpicture}[scale=0.55]
\draw[thin,gray!35] (-2,-1) grid (6,5);
\coordinate (A) at (0,0); \coordinate (B) at (3,0); \coordinate (C) at (1,2);
\coordinate (M) at (5,4); \coordinate (N) at (-1,1);
\draw[thick] (A)--(B)--(C)--cycle;
\draw[-{Stealth},thick,repred] (B)--(M); \draw[-{Stealth},thick,repred] (A)--(N);
\draw[gray,dashed] (N)--(M);
\node[ptcross] at (A) {}; \node[below left] at (A) {\Pt{A}};
\node[ptcross] at (B) {}; \node[below right] at (B) {\Pt{B}};
\node[ptcross] at (C) {}; \node[above left] at (C) {\Pt{C}};
\node[ptcross,repred] at (M) {}; \node[above,repred] at (M) {\Pt{M}};
\node[ptcross,repred] at (N) {}; \node[above,repred] at (N) {\Pt{N}};
\end{tikzpicture}
\end{center}
\textbf{Observation}~: les points $C$, $M$ et $N$ semblent align\'es.
\end{enumerate}
\end{minipage}
\begin{enumerate}[start=2]
\item \textbf{D'une part}, d'apr\`es la relation de Chasles~:
\[\Vj{CM}=\Vj{CB}+\Vj{BM}=\Vj{CB}+2\,\Vj{AC}.\]
\textbf{D'autre part}, d'apr\`es la relation de Chasles~:
\[\Vj{CN}=\Vj{CA}+\Vj{AN}=-\Vj{AC}+\dfrac12\,\Vj{BC}=-\dfrac12\,\Vj{CB}-\Vj{AC}.\]
\item On remarque que $\Vj{CM}=-2\,\Vj{CN}$. En effet~:
\[-2\,\Vj{CN}=-2\left(-\dfrac12\,\Vj{CB}-\Vj{AC}\right)=\Vj{CB}+2\,\Vj{AC}=\Vj{CM}.\]
Donc les vecteurs $\Vj{CM}$ et $\Vj{CN}$ sont colin\'eaires.\\[3pt]
\textbf{Ainsi}~: les points $C$, $M$ et $N$ sont align\'es.
\end{enumerate}
}

\Corr{2}{
\begin{minipage}{\linewidth}
\begin{enumerate}
\item $\Vj{AB}$ correspond \`a 4 carreaux \`a droite, donc $M$ est \`a 1 carreau \`a droite de $A$.\\
$\Vj{BC}$ correspond \`a 2 carreaux \`a droite et 4 en haut, donc $\Vj{BN}=\dfrac34\Vj{BC}$ \`a $1{,}5$ carreau \`a droite et 3 en haut \`a partir de $B$.
\begin{center}
\begin{tikzpicture}[scale=0.55]
\draw[thin,gray!35] (-1,-1) grid (7,5);
\coordinate (A) at (0,0); \coordinate (B) at (4,0); \coordinate (C) at (6,4);
\coordinate (M) at (1,0); \coordinate (N) at (5.5,3);
\draw[thick] (A)--(B)--(C)--cycle;
\draw[repred,thick] (M)--(N);
\node[ptcross] at (A) {}; \node[below left] at (A) {\Pt{A}};
\node[ptcross] at (B) {}; \node[below right] at (B) {\Pt{B}};
\node[ptcross] at (C) {}; \node[above] at (C) {\Pt{C}};
\node[ptcross,repred] at (M) {}; \node[below,repred] at (M) {\Pt{M}};
\node[ptcross,repred] at (N) {}; \node[right,repred] at (N) {\Pt{N}};
\end{tikzpicture}
\end{center}
\textbf{Observation}~: les droites $(MN)$ et $(AC)$ semblent parall\`eles.
\end{enumerate}
\end{minipage}
\begin{enumerate}[start=2]
\item \textbf{Tout d'abord}, exprimons $\Vj{MB}$ en fonction de $\Vj{AB}$~:
\[\Vj{MB}=\Vj{MA}+\Vj{AB}=-\dfrac14\,\Vj{AB}+\Vj{AB}=\dfrac34\,\Vj{AB}.\]
\textit{Justifications}~: relation de Chasles~; $\Vj{MA}=-\Vj{AM}=-\dfrac14\Vj{AB}$~; puis $-\dfrac14+1=\dfrac34$.\\[4pt]
\textbf{Ensuite},
\[\Vj{MN}=\Vj{MB}+\Vj{BN}=\dfrac34\,\Vj{AB}+\dfrac34\,\Vj{BC}=\dfrac34\left(\Vj{AB}+\Vj{BC}\right)=\dfrac34\,\Vj{AC}.\]
\textit{Justifications}~: relation de Chasles~; on remplace $\Vj{MB}$ et $\Vj{BN}$ par leurs expressions~; on factorise par $\dfrac34$~; relation de Chasles.\\[4pt]
\textbf{Conclusion}~: $\Vj{MN}=\dfrac34\,\Vj{AC}$, donc les vecteurs $\Vj{MN}$ et $\Vj{AC}$ sont colin\'eaires.\\[3pt]
\textbf{Ainsi}~: les droites $(MN)$ et $(AC)$ sont parall\`eles.
\end{enumerate}
}

\Corr{3}{
\begin{enumerate}
\item $\Vj{AM}=-\dfrac12\Vj{AB}$~: $M$ est \`a 2 carreaux \`a gauche et 1 en bas de $A$. $\Vj{AN}=-\dfrac12\Vj{AC}$~: $N$ est \`a 1 carreau \`a droite et 1 en bas de $A$.
\begin{center}
\begin{tikzpicture}[scale=0.55]
\draw[thin,gray!35] (0,0) grid (9,5);
\coordinate (A) at (4,2); \coordinate (B) at (8,4); \coordinate (C) at (2,4);
\coordinate (M) at (2,1); \coordinate (N) at (5,1);
\draw[thick] (A)--(B)--(C)--cycle;
\draw[repred,thick] (M)--(N);
\node[ptcross] at (A) {}; \node[below left] at (A) {\Pt{A}};
\node[ptcross] at (B) {}; \node[below right] at (B) {\Pt{B}};
\node[ptcross] at (C) {}; \node[above left] at (C) {\Pt{C}};
\node[ptcross,repred] at (M) {}; \node[below,repred] at (M) {\Pt{M}};
\node[ptcross,repred] at (N) {}; \node[below,repred] at (N) {\Pt{N}};
\end{tikzpicture}
\end{center}
\item D'apr\`es la relation de Chasles~:
\[\Vj{MN}=\Vj{MA}+\Vj{AN}=\dfrac12\,\Vj{AB}-\dfrac12\,\Vj{AC}=\dfrac12\left(\Vj{AB}-\Vj{AC}\right).\]
Or $\Vj{AB}-\Vj{AC}=\Vj{AB}+\Vj{CA}=\Vj{CA}+\Vj{AB}=\Vj{CB}=-\Vj{BC}$.\\
Donc $\Vj{MN}=-\dfrac12\,\Vj{BC}$.
\item Comme $\Vj{MN}=-\dfrac12\,\Vj{BC}$, les vecteurs $\Vj{MN}$ et $\Vj{BC}$ sont colin\'eaires~: les droites $(MN)$ et $(BC)$ sont parall\`eles.
\end{enumerate}
}

\Corr{4}{
\begin{enumerate}
\item Figure (les points $A$, $E$ et $F$ semblent align\'es)~:
\begin{center}
\begin{tikzpicture}[scale=0.75]
\coordinate (A) at (0,0); \coordinate (B) at (3,0); \coordinate (D) at (1,2); \coordinate (C) at (4,2);
\coordinate (E) at ($(B)!0.5!(C)$); \coordinate (F) at ($(D)!2!(C)$);
\draw (A)--(B)--(C)--(D)--cycle;
\draw[repred] (C)--(F); \draw[repred,dashed] (A)--(F);
\node[below] at (A) {\Pt{A}}; \node[below] at (B) {\Pt{B}}; \node[above] at (C) {\Pt{C}}; \node[above] at (D) {\Pt{D}};
\node[ptcross,repred] at (E) {}; \node[right,repred,inner sep=2pt] at (E) {\Pt{E}};
\node[ptcross,repred] at (F) {}; \node[above,repred] at (F) {\Pt{F}};
\end{tikzpicture}
\end{center}
\item Comme $ABCD$ est un parall\'elogramme, $\Vj{BC}=\Vj{AD}$. Donc
\[\Vj{AE}=\Vj{AB}+\Vj{BE}=\Vj{AB}+\dfrac12\,\Vj{BC}=\Vj{AB}+\dfrac12\,\Vj{AD}.\]
\item De m\^eme $\Vj{DC}=\Vj{AB}$, donc
\[\Vj{AF}=\Vj{AD}+\Vj{DF}=\Vj{AD}+2\,\Vj{DC}=\Vj{AD}+2\,\Vj{AB}.\]
\item On remarque que $2\,\Vj{AE}=2\,\Vj{AB}+\Vj{AD}=\Vj{AF}$. Les vecteurs $\Vj{AE}$ et $\Vj{AF}$ sont donc colin\'eaires~: les points $A$, $E$ et $F$ sont align\'es.
\end{enumerate}
}

\Objectif{7}{Caract\'erisation vectorielle du milieu d'un segment}

\Corr{5}{
\begin{enumerate}
\item Figure~:
\begin{center}
\begin{tikzpicture}[scale=0.5]
\coordinate (A) at (0,0); \coordinate (B) at (3,0); \coordinate (D) at (1,2); \coordinate (C) at (4,2);
\coordinate (E) at ($(A)!2!(B)$); \coordinate (F) at ($(A)!2!(D)$);
\draw (A)--(B)--(C)--(D)--cycle;
\draw[repred] (B)--(E); \draw[repred] (D)--(F); \draw[repred,dashed] (E)--(F);
\node[below] at (A) {\Pt{A}}; \node[below] at (B) {\Pt{B}}; \node[above right,inner sep=1pt] at (C) {\Pt{C}}; \node[left] at (D) {\Pt{D}};
\node[ptcross,repred] at (E) {}; \node[below,repred] at (E) {\Pt{E}};
\node[ptcross,repred] at (F) {}; \node[above,repred] at (F) {\Pt{F}};
\end{tikzpicture}
\end{center}
\item $B$ est le milieu de $[AE]$, donc $\Vj{BE}=\Vj{AB}$. Comme $ABCD$ est un parall\'elogramme, $\Vj{CB}=\Vj{DA}=-\Vj{AD}$. Ainsi
\[\Vj{CE}=\Vj{CB}+\Vj{BE}=-\Vj{AD}+\Vj{AB}=\Vj{AB}-\Vj{AD}.\]
\item De m\^eme, $D$ est le milieu de $[AF]$, donc $\Vj{DF}=\Vj{AD}$, et $\Vj{CD}=\Vj{BA}=-\Vj{AB}$. Ainsi
\[\Vj{CF}=\Vj{CD}+\Vj{DF}=-\Vj{AB}+\Vj{AD}.\]
\item On constate que $\Vj{CF}=-\Vj{CE}$, c'est-\`a-dire $\Vj{CE}+\Vj{CF}=\vec 0$~: c'est la caract\'erisation vectorielle du milieu. Donc $C$ est le milieu du segment $[EF]$.
\end{enumerate}
}

\end{multicols}
\end{document}
